\scanpage{54}
Antipodisme
\deleted{\diagram{\begin{tikzcd}[ampersand replacement=\&]
x\arrow[r]\&2o\Pi(x)-x
\end{tikzcd}}}
\[
\uline a=-\mathrm{id}+\deleted{\unsure{2o}}
2\Pi\otimes o=-\mathrm{id}+\Pi\otimes o'
\]
\[
-\uline a=\mathrm{id}-\Pi\otimes o'
\]
\uline{NB}\quad$\uline a$ défini sss $o'$ défini i.e. sss $\rho'$ inversible
(c'est le cas si $f_V$ lisse). Notons que
\[
\langle o',\Pi\rangle=2,
\]
donc $-\uline a$ est une réflexion,
\inserted{sauf si en certains pts on a car.$=2$, \uline{et}
\deleted{\unsure{dans}}\deleted{\unsure{ce qui}}\ \unsure{rang impair}}.
$\uline a$ commute aux $\tau_i$, il invarie $o'$, $o$ etc.

\margin{\uline{NB} Si $n$ pair ou $2$ inv., alors:
$\uline a$ défini \uline{et} est une réflexion \uuline{str. orth.}
$\Longleftrightarrow\sigma'\rho'$ inv., i.e. $\sigma'$ et $\rho'$ inv.}
\margin{\uline{Attention}, \unsure{il se peut que} $\unsure{c|o}$ n'existe pas à dist. finie,
mais $\unsure{c'|o'}$ existe\ldots
\[
\unsure{\text{si }n\text{ impair}}\qquad
\rho'o'=\widetilde o\quad\unsure{\text{pour }2\neq0}
\]
\unsure{donc $o'$ pour $2\neq0$}, donc
$\unsure{\rho o'=\rho\widetilde o}$.
\unsure{m conditions pour les parties si $2$ inv.}
Mais en car. $2$, pour $n$ \unsure{impair}, on a $o'=0$
d'où $\uline a=\mathrm{id}$.}
\figure{\begin{tikzpicture}
\draw (0,3)..controls(.2,2)and(2,2.6)..(1.6,1.8)..controls(1.3,1)and(.2,.5)..(0,0);
\draw (4,3)..controls(3.6,2.7)and(3.6,2.6)..(4.2,2.4)--(4.2,.4)..controls(4.2,.1)and(4.4,0)..(4.8,0);
\end{tikzpicture}}
\[
\uline a(s_0)=\uline a(a_{-1})
=\deleted{\unsure{o}}o'-s_0=\widetilde s_0
\]
\uline{NB}\quad Dans $\mathcal E$, l'ens. des pts fixes de $\uline a$ est
\deleted{\unsure{\text{la droite}}}
la droite engendrée par $o'$ si $2$ inv.
--- si $2\,1_S=0$ i.e. $S$ de car. $2$, alors l'ens des pts fixes est $V$.
Dans $P(\mathcal E)$, l'ens des pts fixes est $P(V)\cup\{o\}$
--- qui se réduit à $P(V)$ en car. $2$.

Donc $s_0\neq\widetilde s_0=\uline a s_0$ dans $P(\mathcal E)=X$.

Considérons le \deleted{droite} plan engendré par $s_0$ et $\widetilde o$
\deleted{\unsure{(et $o'$)}}
\[
\begin{aligned}
g(\lambda,\mu)&=f(\lambda s_0+\mu\widetilde o)
=\deleted{\unsure{\lambda^2}}
\lambda\mu\underbrace{\varphi(s_0,\widetilde o)}_{-\zeta_0=-\sigma}
+\mu^2f(\widetilde o)\\
&=-\sigma\lambda\mu-\sigma'\rho'\mu^2
=-\mu(\sigma\lambda+\sigma'\rho'\mu)
\end{aligned}
\]
1) \uline{$n$ pair}: $\sigma=\sigma'$, $\rho=2\rho'$
\[
f(\lambda s_0+\mu\widetilde o)=-\mu\sigma(\lambda+\rho'\mu)
\]
c'est $=0$ sss $\sigma=0$.

c'est $\neq0$ en chaque fibre $\Longleftrightarrow\sigma$ inv. i.e. $\sigma'$ inv.
\[
\Longleftrightarrow\quad
\deleted{\unsure{\text{donc on a}}}
Q_X\cap D_0\ \text{est un rev. étale de longueur }2\text{ de }S,
\]
dont $s_0$ est une section,
[$\uline a_Xs_0$ est la section complémentaire,
quand $\uline a_X$ est définie]
\[
\Longleftrightarrow Q_X\unsure{\not\supset|\neq}D_0\ \text{en tt fibre}
\]
(C'est le cas si $f_{\mathcal E}$ lisse)
\[
\Longleftrightarrow\Psi(\widetilde o)\neq0\ \text{chaque fibre}
\Longleftrightarrow(\widetilde o)^{\unsure{\perp}}=V
\quad\text{chaque fibre}
\]
\margin{\uline{NB} Si $\sigma$ inv. \deleted{\unsure{\text{précisons}}}
\[
f(\lambda s_0+\mu\widetilde o)=0,\quad
\deleted{\unsure{\mu}}\inserted{\mu}\text{ inv.}
\]
sss $\lambda+\rho'\mu=0$ i.e. $\lambda=-\rho'\mu$
i.e. $(\lambda,\mu)\parallel(-\rho',1)$
donc $\lambda s_0+\mu\widetilde o$ multiple de $-\rho's_0+\widetilde o$
ou encore (si $\rho'$ inv.) de
$-s_0+\dfrac{\widetilde o}{\rho'}=-s_0+o'=\widetilde s_0$.}
2) \uline{$n$ impair}\quad $\sigma=2\sigma'$, $\rho=\rho'$,
\[
f(\lambda s_0+\mu\widetilde o)=-\mu\sigma'(2\lambda+\rho'\mu)
\]
\mspage{53}

